\documentclass[11pt]{article}
\usepackage[T1]{fontenc}
\usepackage{lmodern,amsmath,amssymb,amsthm,mathtools}
\usepackage[margin=1in]{geometry}
\usepackage{booktabs,longtable,array,microtype,enumitem}
\usepackage{xcolor,xurl,hyperref,fancyhdr,listings}
\hypersetup{colorlinks=true,linkcolor=blue!45!black,urlcolor=blue!45!black,citecolor=blue!45!black,
pdftitle={Congruence, Radical Norms, and Direct Proof Certificates for Diophantine ZFC Consistency},
pdfauthor={AI-assisted research contribution prepared for Vladimir Reshetnikov},
pdfsubject={Explicit Diophantine equations, ordinary proofs, and reproducible finite verification}}
\lstset{basicstyle=\ttfamily\small,breaklines=true,columns=fullflexible,keepspaces=true,frame=single}
\setcounter{tocdepth}{1}
\newtheorem{theorem}{Theorem}[section]
\newtheorem{lemma}[theorem]{Lemma}
\newtheorem{proposition}[theorem]{Proposition}
\newtheorem{corollary}[theorem]{Corollary}
\theoremstyle{remark}\newtheorem{remark}[theorem]{Remark}
\newcommand{\Z}{\mathbb Z}\newcommand{\N}{\mathbb N}\newcommand{\Q}{\mathbb Q}
\newcommand{\Con}{\operatorname{Con}}\newcommand{\Pair}{\operatorname{Pair}}
\newcommand{\F}{\mathcal F}\newcommand{\vp}{\nu_3}
\newcommand{\repo}{\nolinkurl{a88aefd7ddab4771534bb049c67350e397c91b9a}}
\pagestyle{fancy}\fancyhf{}
\fancyhead[L]{\small Diophantine consistency: arithmetic reductions}
\fancyhead[R]{\small ProveIt research contribution}
\fancyfoot[C]{\thepage}\setlength{\headheight}{14pt}
\title{Congruence, Radical Norms, and Direct Proof Certificates\\for Diophantine ZFC Consistency}
\author{Research contribution prepared for Vladimir Reshetnikov\\\small AI-assisted mathematical research and computational verification}
\date{24 October 2026}
\begin{document}
\maketitle
\begin{abstract}
We supply explicit polynomial equations representing syntactic inconsistency
of ZFC by simplifying two different existing constructions. A family over
unrestricted integers combines congruence recovery, fixed-prime separation
of radicals, and products of integer factors. Its principal member has
16 variables, exact degree 124, and 118 arithmetic operations, improving
the previous 16-variable degree-184, 135-operation point in both measures.
A separate direct proof-certificate compiler produces smaller explicit
quartics with modest numerical coefficients, using neither K215 nor a
Jones universal polynomial. We also remove a congruence restriction from
two positive-variable equations and prove a product-certificate lemma
that computes its own coefficient bound.

Every numerical claim concerns a supplied finite polynomial or arithmetic
circuit. Ordinary proofs of the new reductions are distinguished from
finite checks and inherited calibration theorems. There is no claim of a
globally simplest equation, a new minimum-variable record, or a completed
proof-assistant formalization of the entire ZFC equivalence.
\end{abstract}
\tableofcontents\newpage

\section{Results, conventions, and the exact objective}
The target is a fixed polynomial with integer coefficients for which
\[
 \mathrm{ZFC}\vdash
 \bigl(\neg\exists\boldsymbol{x}\in D^m\;P(\boldsymbol{x})=0\bigr)
 \longleftrightarrow\Con(\mathrm{ZFC}),
\]
with an explicitly stated domain $D$. Consistency means absence of a
finite syntactic derivation of a contradiction in the calibrated calculus.
No consistency or inconsistency assumption is used to choose a coefficient.
Natural numbers include zero; positive integers do not.

The comparison snapshot is \repo. We inspected the active research and
recent incoming archives, rather than relying on the older summary of
packages 31--75. The repository already contains nine-positive and
ten-integer equations. Our modest-degree integer equations do not improve
those variable counts. The direct route has much smaller numerical
coefficients but many more variables. Several measures must be retained.
\begin{table}[ht]
\centering\small
\begin{tabular}{@{}lrrrl@{}}
\toprule
Domain & Variables & Degree & Operations & Comparison\\
\midrule
$\Z$ & 13 & 1,072 & 195 & previous 1,088 / 195\\
$\Z$ & 14 & 414 & 148 & previous 696 / 150\\
$\Z$ & 15 & 218 & 130 & previous 214 / 139\\
$\Z$ & \textbf{16} & \textbf{124} & \textbf{118} & previous 184 / 135\\
$\Z_{>0}$ & 12 & 70 & 98 & previous 70 / 100\\
$\Z_{>0}$ & 13 & 51 & 88 & previous 51 / 90\\
\bottomrule
\end{tabular}
\caption{New literal native-family polynomials. Previous entries are selected
relevant points in the inspected repository, not assertions of global optima.}
\end{table}
The 16-integer result lowers degree by $32.6\%$ and operations by $12.6\%$
against the degree-184, 135-operation parent. It also improves the separate
degree-192, 134-operation point in both measures. The 14-integer equation
has a lower degree than both the 696/150 root-fusion equation and the
652/230 fixed-prime equation. The 792/137 equation remains a lower-cost
alternative. The 15-integer entry is a degree--cost tradeoff.

\begin{table}[ht]
\centering\small
\begin{tabular}{@{}lrrrrr@{}}
\toprule
Direct quartic over $\N$ & Variables & Rows & Terms & Height & Degree\\
\midrule
Parent, affine elimination & 3,614 & 2,954 & 15,996 & 127,778,545 & 4\\
New, affine elimination & \textbf{2,473} & 2,089 & 20,580 & 127,808,289 & 4\\
Parent, retained affine rows & 6,183 & 5,523 & 17,430 & 1,993 & 4\\
New, retained affine rows & \textbf{4,860} & 4,476 & 14,719 & \textbf{1,837} & 4\\
\bottomrule
\end{tabular}
\caption{Fully expanded direct polynomials. Height is the maximum absolute
coefficient. Rows are quadratic residuals before summing their squares.}
\end{table}
The first new direct equation has $31.6\%$ fewer variables than its parent,
at the cost of more expanded monomials and a slightly larger coefficient
height. The retained-row variant reduces the variable, row, and monomial
counts and the coefficient height against its corresponding parent,
while retaining degree four. Its coefficients need at most eleven magnitude bits.



An operation is one live binary addition, subtraction, or multiplication.
Squaring and multiplication by a fixed numeral each count as one operation.
Input and fixed-numeral leaves are free. Polynomial abbreviations are
shared registers, not additional existential variables. For the direct
quartics we give literal variable, residual, monomial, and coefficient
counts; these do not substitute for arithmetic-circuit counts.

\subsection{New results and imported calibration}
The new algebraic reductions have ordinary proofs below. They identify
solutions with explicitly stated parent systems. The final interpretation
as refutability uses two inherited ordinary calibrations:
\begin{enumerate}[leftmargin=*]
\item The native branch uses the signed native/Pell compiler and literal
K215 blank-run calibration, with the complete source-34 or compact44
coefficient tuple appropriate to that branch.
\item The direct branch uses the C222 proof calculus, Horn certificates,
bounded-sequence coding, and the ordinary C222--ZF--ZFC finite-refutation
translation. It uses no native tuple or K215 simulation.
\end{enumerate}
These theorems, their provenance, and their proof boundaries are retained
in the package. The new arguments use finite arithmetic, induction,
polynomial identities, Pell recurrences, and standard three-square
arithmetic, all formalizable in ZFC. Composing ordinary ZFC proofs gives
the target theorem. A new Lean/Rocq derivation, or a literal derivation in
the repository's encoded object-language ZFC, is not supplied. Numerical
checks do not independently establish the unbounded calibration theorems.

\section{Three elementary tools}
\subsection{A product of integer factors}
\begin{lemma}[Product finalizer]\label{lem:units}
Let integer-valued polynomials $E_1,\ldots,E_s$ satisfy
$E_i(\boldsymbol{x})\not\equiv2\pmod4$ on every integer assignment. Then
\[
 \prod_{i=1}^s(1+E_i)-1=0
 \quad\Longleftrightarrow\quad E_1=\cdots=E_s=0.
\]
\end{lemma}
\begin{proof}
An integer product equal to one has each factor equal to $1$ or $-1$.
The latter would give $E_i=-2\equiv2\pmod4$. Every factor must therefore
equal one. The converse is immediate.
\end{proof}
The residue hypothesis is essential: two arbitrary residuals equal to
$-2$ would otherwise create a spurious zero.

\subsection{A fixed prime separates square roots}
\begin{lemma}[Valuation separation]\label{lem:valuation}
Let $A_1,\ldots,A_s\in\Q^\times$ have pairwise distinct even $3$-adic
valuations. If a signed sum of square roots of the $A_i$ is rational,
every one of those roots is rational. Negative radicands are allowed.
\end{lemma}
\begin{proof}
Group the radicands by their classes in $\Q^\times/(\Q^\times)^2$.
Write roots in each class as $r_i\sqrt d$, with rational nonzero $r_i$.
Distinct square classes give linearly independent radical monomials in
the associated multiquadratic extension. In each nontrivial class the
signed coefficient sum must vanish. But its rational coefficients have
pairwise distinct valuations, half the differences of the original
valuations. A sum with a unique term of least valuation cannot vanish.
There is therefore no nontrivial class. This also applies to negative $d$.
\end{proof}

\subsection{Positivity with a fixed residue}
Define $\Pair(x,y,t)=xy-(t^2-x-y+1)^2$. For integers $x,y$, some integer
$t$ makes this positive exactly when $x,y>0$. Positivity implies $xy>0$;
if both were negative, the subtracted square would exceed $xy$.
Conversely, set $L=x+y-1$. The nearest square to $L$ has squared error
strictly below $L\le xy$ unless $L$ itself is square, when the error is zero.

\begin{lemma}[Pentagonal positivity]\label{lem:positive}
For $N\in\Z$ put
\[
 S(N,u_1,u_2)=24N-21-(6u_1-1)^2-(6u_2-1)^2.
\]
Then $N>0$ if and only if some integers $u_1,u_2,g$ satisfy $S=g^2$.
Always $S\equiv1\pmod3$ and $S\equiv1\pmod4$.
\end{lemma}
\begin{proof}
The equation gives $24N-21\ge2$, hence $N\ge1$. Conversely, the
classical theorem on three generalized pentagonal numbers gives
\[
 24N-21=(6u_1-1)^2+(6u_2-1)^2+(6u_3-1)^2.
\]
It follows from the three-square theorem by an elementary norm-$9$
lifting. Remove the largest power $9^e$ from $24N-21$, represent the
remaining positive number, which is $3\pmod8$, as three odd squares,
and choose signs before each application of
\[
 L=\begin{pmatrix}1&2&2\\2&1&-2\\-2&2&-1\end{pmatrix},
 \qquad L^{\mathsf T}L=9I.
\]
The signs can make $x-y-z\not\equiv0\pmod3$ unless all three
coordinates vanish modulo three. The resulting coordinates are all
nonzero modulo three. This restores the removed powers of nine while
preserving oddness. Choose final signs to make each coordinate $-1\pmod6$.
This classical argument is discussed by Sun and Oh; neither it nor the
matrix identity is new\cite{sunpent,oh}. The residues follow directly.
\end{proof}

\section{An explicit family over unrestricted integers}\label{sec:integer}
\subsection{Fixed coefficients and base coordinates}
The six fixed coefficients are the coherent source-34 tuple
\[
 \beta=2^{1064921}-1,\quad K_0,\quad M_C,\quad M_F,
 \quad A_c=2^{1113},\quad B_c=2^{666104}.
\]
Here $M_F$ already includes the native shift by $\beta$. The complete
minimal little-endian numeral payloads and their binding file are supplied.
They are not unspecified parameters. The retained finite tile recipe and
source pin explain their mathematical origin; no choice based on halting
is involved. Their uncompressed total is 615,883 bytes.

Use the thirteen base coordinates
\[
 j,C,z,h,r,v,w,T,c,n,t,u_1,u_2\in\Z.
\]
Uppercase $C$ and lowercase $c$ are different coordinates. Set
\begin{align}
 J&=6j+1,&q&=\beta J+1,&W&=12w+4,\nonumber\\
 X&=qW,&Y&=qv,&E&=XY,\nonumber\\
 a&=(X+1)Y,&H&=4a+3,&\Delta&=(a+2)^2-1,\label{eq:front}\\
 F&=(K_0+W)C-z(q-1),\nonumber\\
 R&=[q(q-F)-C](q^2-1)+(M_C+qM_F)J,\nonumber\\
 k&=R+1+hE,&\eta&=c-Yk,&\zeta&=k-\eta,\label{eq:rank}\\
 I&=(2r+1)c^2,&D_f&=1+\Delta I^2,\nonumber\\
 D_\tau&=1+XY^2(XY^2+1)k^2,&D_m&=1+\Delta c^2,&b&=X+ac.
 \label{eq:radicands}
\end{align}
Define the determinant and positivity polynomials
\begin{align}
 G&=(A_cF)(B_cC)-(q-A_cF-B_cC)^2,\nonumber\\
 N&=G\eta\zeta-(t^2-G-\eta\zeta+1)^2,\nonumber\\
 S&=24N-21-(6u_1-1)^2-(6u_2-1)^2.\label{eq:guard}
\end{align}
The carrier scales are $M=2q^2$ and $U=MH$.
For the thirteen-variable equation put
\[
 Q=\Delta^2 I^2,\qquad V=D_f(cT-R)-c.
\]
For equations with a retained coordinate $f$, put instead
\begin{equation}
 Q=\Delta(f^2-1),\qquad V=f^2(cT-R)-c.\label{eq:retainedf}
\end{equation}
These agree when $f^2=D_f$. Finally set
\begin{equation}
 \delta=Q-1,\qquad C_y=QV^2-1,\qquad
 Z=U(n+V)+Mb+H(v-q^2).\label{eq:carrier}
\end{equation}

\subsection{A finite polynomial definition of the radical merger}
For $s$ radicands define the universal polynomial
\[
 \F_{2^s}(x;A_1,\ldots,A_s)=
 \prod_{\varepsilon\in\{-1,1\}^s}
 \left(x-\sum_{i=1}^s\varepsilon_i\sqrt{A_i}\right).
\]
Changing any radical sign permutes the factors, so this lies in
$\Z[x,A_1,\ldots,A_s]$. Each radicand has exponent at most $2^{s-1}$.
Write
\[
 \mathcal C_s(x;A_1,\ldots,A_{s-1};B,\delta)
 =\delta^{2^{s-1}}\F_{2^s}(x;A_1,\ldots,A_{s-1},B/\delta).
\]
The denominator is canceled symbolically. This is a finite polynomial
definition, including when $\delta=0$. The literal JSON circuits use
only addition, subtraction, and multiplication.

With $B_y=U^2C_y$, specialize
\begin{align*}
 N_5&=\mathcal C_5(Z;M^2D_m,9U^2D_f,81U^2D_\tau,729U^2S;B_y,\delta),\\
 N_4&=\mathcal C_4(Z;M^2D_m,9U^2D_\tau,81U^2S;B_y,\delta),\\
 N_3&=\mathcal C_3(Z;M^2D_m,9U^2D_\tau;B_y,\delta),\\
 N_2&=\bigl[\delta(Z^2-M^2D_m)-B_y\bigr]^2-4\delta M^2D_mB_y.
\end{align*}
The weights change after retaining the strong root: there is no reason
to preserve a vacant valuation position. This lowers the literal cost.
Put
\begin{align*}
 F_f&=f^2-\Delta I^2=1+(f^2-D_f),\\
 F_g&=g^2+22+(6u_1-1)^2+(6u_2-1)^2-24N=1+(g^2-S),\\
 F_\tau&=\tau^2-XY^2(XY^2+1)k^2=1+(\tau^2-D_\tau).
\end{align*}
The four equations are
\begin{align}
 P_{13}&=N_5=0,\label{eq:p13}\\
 P_{14}&=(1+N_4)F_f-1=0,\label{eq:p14}\\
 P_{15}&=(1+N_3)F_fF_g-1=0,\label{eq:p15}\\
 \boxed{P_{16}}&\boxed{=(1+N_2)F_fF_gF_\tau-1=0.}\label{eq:p16}
\end{align}
The added coordinates are none, $f$, $(f,g)$, and $(f,g,\tau)$,
respectively. All coordinates are unrestricted integers. The last line,
together with the polynomial abbreviations above, is the recommended
modest-degree equation in this report.

\subsection{Recovering retained equations by residues}
The affine formulas imply that $q$ is even, $q\equiv2\pmod3$, and
$W$ is a nonzero multiple of four with $W\equiv1\pmod3$. Thus $X$ is
divisible by eight, $Y$ and $a$ are even, and $\Delta\equiv3\pmod4$.
The integer $D_f$ is $0$ or $1\pmod4$, while $D_\tau\equiv1\pmod4$.
Lemma~\ref{lem:positive} gives $S\equiv1\pmod4$. Therefore
\[
 F_f\equiv f^2+I^2\in\{0,1,2\},\quad
 F_g\equiv g^2\in\{0,1\},\quad F_\tau\equiv\tau^2\in\{0,1\}\pmod4.
\]
None can equal $-1$.

Both $M$ and $U$ are divisible by eight. Modulo four every norm monomial
using a radicand vanishes, and
\begin{equation}
 N_s\equiv\delta^{2^{s-1}}Z^{2^s}\in\{0,1\}\pmod4\qquad(s\ge2).
 \label{eq:normmod4}
\end{equation}
There is no assumption that $\delta$ is odd; it need not be.
Lemma~\ref{lem:units} makes each product equation exactly its displayed
simultaneous system. In particular $f^2=D_f$ is recovered before using
\eqref{eq:retainedf} as the normalized strong Pell expression.

\subsection{Square extraction before positivity}
The order of the proof matters. Since $J\ne0$ and $\beta$ is large,
$|q|\ge16$ and $|X|\ge64$. If $v=0$, then $a=0$ and $\Delta=3$.
Otherwise $|Y|\ge16$ and $|a|=|X+1||Y|\ge1008$, so $\Delta>1$.
Always $H\ne0$ and $\Delta\ge3$.

After recovery of retained factors, if any, $Q=\Delta^2I^2$.
For $c=0$, $Q=0$ and $\delta=-1$; for $c\ne0$, $Q>1$ and
$\delta>0$. Thus the cleared norm can be interpreted as the product of
conjugate linear factors. Alternatively, the universal congruence
$\delta\equiv-1\pmod3$ already excludes a zero denominator.

Modulo three we have $X\equiv2$, $a=(X+1)Y\equiv0$, and
\[
 H\equiv\Delta\equiv0,\quad
 D_m\equiv D_f\equiv D_\tau\equiv S\equiv1,\quad
 Q\equiv0,\quad\delta\equiv C_y\equiv-1.
\]
For $D_\tau$, use $XY^2(XY^2+1)\equiv0\pmod3$ for either possible
residue of $Y^2$. Consequently $M$ is a unit at three and
$e=\vp(U)=\vp(H)\ge1$. In the five-root case the valuations of the
scaled radicands, in the order main, auxiliary, strong, first, guard, are
\[
 0,\quad2e,\quad2e+2,\quad2e+4,\quad2e+6.
\]
In the retained-strong family they are $0,2e,2e+2,2e+4$, omitting
any further retained roots. These are distinct even integers.
The norm zero gives a rational signed sum. Lemma~\ref{lem:valuation}
makes each unretained root rational, hence integer when its radicand is
integer. The guard acquires a root in this way or has a retained root.
Lemma~\ref{lem:positive} and the pair guard then give
\begin{equation}
 N>0,\qquad G>0,\qquad\eta\zeta>0.\label{eq:recoveredguard}
\end{equation}
If $c=0$, then $\eta\zeta=-Y(Y+1)k^2\le0$, a contradiction.
Thus $c\ne0$ and $Q>1$ at every solution.

\begin{lemma}[Adjacent-coefficient rational descent]\label{lem:descent}
If $Q>1$ and $V$ are integers and $y\in\Q$ satisfies
$QV^2-(Q-1)y^2=1$, then $y\in\Z$.
\end{lemma}
\begin{proof}
Write $y=A/d$ in lowest terms; then $d^2\mid Q-1$.
For $|V|=1$ the claim is immediate. After changing signs if needed, apply
\[
 V'=(2Q-1)V-2(Q-1)y,\qquad y'=(2Q-1)y-2QV.
\]
The first coordinate is integral, the norm is unchanged, and the reduced
denominator of the second remains $d$, since $2Q-1\equiv1\pmod d$.
For $V>1$, the norm gives $y>V$ and $(Q-1)y<QV$, whence
$-V<V'<V$; it also excludes $V'=0$. Descent on $|V|$ reaches one
without changing the denominator.
\end{proof}
The auxiliary radical has square $C_y/\delta$ and supplies the rational
$y$ in this lemma. It is therefore integral too.

\subsection{The two congruences and their inverse map}
The vanishing factor has the form
\begin{equation}
 Z=Md+UL,\qquad d^2=D_m,\quad L\in\Z,\label{eq:linearcarrier}
\end{equation}
where signs have been absorbed into the recovered roots. The integer
$L$ is their appropriate weighted sum. Because
$H=4(X+1)qv+3\equiv3\pmod q$, $H$ is odd, and $3\nmid q$,
\[
 \gcd(H,M)=\gcd(H,2q^2)=1.
\]
Modulo $M$, equation~\eqref{eq:linearcarrier} gives
\[
 v=q^2(2v_{\rm old}+1),\qquad
 v_{\rm old}=\frac{v-q^2}{2q^2}\in\Z.
\]
Thus $v\ne0$ and $Y=(2v_{\rm old}+1)q^3$, with exactly the required
odd quotient. Modulo $H$ it gives
$\gamma=(d-b)/H\in\Z$. Neither conclusion was assumed in the
valuation argument. Conversely choose
\begin{equation}
 n=\gamma+L-V-\frac{v-q^2}{M}\in\Z.\label{eq:choose-n}
\end{equation}
In the retained family $L=y+3\tau+9g$, omitting the terms whose roots
are retained; in the five-root family $L=y+3f+9\tau+27g$.

\subsection{The normalized source and its calibration}
The restored source has the same front registers and
\[
 f^2=D_f,\quad\tau^2=D_\tau,\quad QV^2-(Q-1)y^2=1,\quad
 (b+H\gamma)^2=D_m,\quad S=g^2,
\]
with $Y=(2v_{\rm old}+1)q^3$. The transport row holds identically
by the definition of $F$. Put $i=2r+1$ and restore older coordinates by
\begin{align*}
 I_{\rm old}&=\Delta i c^2,&r_{\rm old}&=(\Delta i-1)/2,\\
 o_{\rm old}&=f(cT-R),&j_{\rm old}&=D_fT-1-R\Delta i^2c^3.
\end{align*}
The first quotient is integral since $\Delta i$ is odd. Exact identities
give $I_{\rm old}^2=\Delta(f^2-1)$ and
$fo_{\rm old}-c=cj_{\rm old}-R=V$.
To restore the older ternary positivity row, use $N>0$ to choose
integers $b_1,b_2,b_3$ with
$N=1+b_1^2+b_2^2+b_3(b_3+1)$, by the three-square theorem.
These are reselected coordinates of the inherited source witness, not
additional variables in the new polynomial.

For completeness, the imported canonical native witness has
$q=2^e$ with $e$ odd, $R\equiv3\pmod4$, $X=2^R$, and an odd positive
quotient $Y/q^3$. Its native repunit satisfies $J\equiv1\pmod6$, and
$X/(4q)\equiv1\pmod3$. Thus the affine choices $J=6j+1$ and
$W=12w+4$ preserve canonical existence. These are properties of the
same complete source-34 coefficient tuple, not assumptions on arbitrary
integer assignments. Furthermore $c=\psi_A(R)$ with $A=a+2$ even
and $c$ odd. Choose
$I_*=\psi_A(Rc)$ and $f=\chi_A(Rc)$.
Pell composition, or expansion of $(D+c\sqrt\Delta)^c$, gives
$c^2\mid I_*$. The recurrence gives $I_*$ odd; hence $i=I_*/c^2$ is odd.
The retained squared-modulus auxiliary lemma, using the established
phase $R\equiv3\pmod4$, gives $V,y$ with
$V\equiv-c\pmod{f^2}$ and $V\equiv-R\pmod c$. Since
$f^2\equiv1\pmod c$, it supplies the integer
\[
 T=\frac{V+c+Rf^2}{cf^2}.
\]
The determinant and strip guards hold for that same canonical witness.
Choose $t$ by the nearest-square argument, the pentagonal representation
of $N$, and then $n$ by \eqref{eq:choose-n}.
This is the same calibrated source family as the retained root-fusion and
fixed-prime proposals, followed by the new algebraic finalizer.

\begin{theorem}[Integer family]\label{thm:integer}
Each polynomial \eqref{eq:p13}--\eqref{eq:p16}, with its stated fixed
tuple and domain, has a zero exactly when the normalized native source
has a witness. Composed with the inherited ordinary native/K215
calibration, absence of zeros is provably equivalent in ZFC to
$\Con(\mathrm{ZFC})$.
\end{theorem}
\begin{proof}
The product finalizer recovers retained equations; valuation separation
recovers other roots; descent recovers the auxiliary integer; the coprime
carrier recovers the cubic scale and main congruence. The displayed
restorations give the source. The preceding canonical construction proves
the reverse implication. Compose the named ordinary calibration and
negate. The imported compiler theorem is not replaced by numerical tests.
\end{proof}

\section{Exact degrees and arithmetic cost}\label{sec:degrees}
The four integer forms have exact degrees, not merely syntactic upper
bounds. Every fixed numeral has degree zero. For the retained-$f$ family,
the relevant register degrees are
\[
\begin{array}{c|rrrrrrrrrrrr}
\text{register}&a&H&\Delta&R&k&Q&V&C_y&U&Z&S&D_\tau\\\hline
\deg&4&4&8&5&5&10&7&24&6&13&28&22.
\end{array}
\]
The weighted radicands in $N_4$ have degrees $14,34,40$, with auxiliary
numerator degree $36$ and denominator degree $10$.

Write a subscript zero for the leading homogeneous part. Useful expressions are
\begin{align*}
q_0&=6\beta j,&W_0&=12w,&Y_0&=q_0v,&X_0&=q_0W_0,\\
a_0&=X_0Y_0,&i_0&=2r,&R_0&=-q_0^3(W_0C-q_0z),&k_0&=R_0+ha_0,\\
U_0&=8q_0^2a_0,&M_0&=2q_0^2,&Q_0&=a_0^2f^2,&V_0&=-f^2R_0,\\
S_0&=-24Y_0^4k_0^4,&(D_\tau)_0&=X_0^2Y_0^4k_0^2.&&&
\end{align*}

\subsection{Four roots and three roots}
In the four-root norm the guard root has degree 20, strictly exceeding
the degrees 17, 7, 13 of the other roots and degree 13 of the center.
Its leading term is therefore
\[
 (N_4)_{400}=Q_0^8(81U_0^2S_0)^8\ne0.
\]
Multiplication by $F_f$, of degree 14 with leading part
$-a_0^2i_0^2c^4$, gives exact degree 414. A coefficient certificate is
\[
 [j^{196}w^{66}v^{98}h^{32}r^2f^{16}c^4]P_{14}
 =-2^{402}3^{302}\beta^{196}\ne0.
\]
For three roots, the first Pell root is uniquely largest, with degree 17.
Hence $N_3$ has degree $8\cdot17+4\cdot10=176$ and leading part
\[
 Q_0^4(9U_0^2(D_\tau)_0)^4.
\]
The factors $F_f,F_g$ add degrees 14 and 28. Thus $P_{15}$ has degree
218, with the explicit coefficient
\[
 [j^{104}w^{38}v^{50}r^2f^8h^{12}c^4]P_{15}
 =-2^{209}3^{151}\beta^{104}\ne0.
\]
These certificates remain nonzero on the actual fixed numerical tuple.

\subsection{The two-root equation and degree 124}
The center $Z$ contains the leading approximation to the auxiliary root.
Write
\[
 Z=UV+s,\qquad s=Un+Mb+H(v-q^2),\qquad\deg s=7.
\]
Then the exact ring identity
\begin{equation}
 \delta(Z^2-M^2D_m)-U^2C_y
 =\delta(s^2+2UVs-M^2D_m)-U^2(V^2-1)
 \label{eq:centeridentity}
\end{equation}
lowers the degree of this inner expression to 30. Therefore $N_2$ has
degree at most 60. Its complete degree-60 part is
\[
 4Q_0^2U_0^2V_0^2(s_0^2-A_0),\quad
 s_0=M_0a_0(4n+c),\quad A_0=M_0^2a_0^2c^2,
\]
equivalently
\[
 (N_2)_{60}=2^{13}q_0^8a_0^8f^8R_0^2n(2n+c)\ne0.
\]
The remaining factors have degrees 14, 28, and 22. Consequently
\[
 \deg P_{16}=60+14+28+22=124.
\]
For a single-coefficient certificate,
\[
 [j^{58}w^{18}v^{24}r^2h^6f^8z^2n^2c^4]P_{16}
 =2^{113}3^{77}\beta^{58}>0.
\]
The powers $h^6,z^2,n^2$ isolate the corresponding terms in
$k_0^6,R_0^2,n(2n+c)$ uniquely. This proves exact degree independently
of any algebraic relations among the other five fixed coefficients.

\subsection{Five roots and degree 1072}
Without retaining $f$, the degrees change to
$\deg Q=22$, $\deg V=19$, $\deg C_y=60$, $\deg Z=25$.
The auxiliary radical satisfies
\[
 C_y/\delta=V^2+\frac{V^2-1}{Q-1}.
\]
For a rigorous homogeneous-degree computation, scale all inputs by a
formal parameter and work in its Laurent-series extension. The two
auxiliary roots are $\pm(V+\text{terms of degree }-3)$.
After multiplication by $U$, their correction has degree 3.
Sixteen conjugate factors retain the degree-25 auxiliary term; in the
other sixteen it cancels the center's $UV$, leaving the guard root of
degree 20 as the unique leading term. Clearing denominators contributes
$16\cdot22$. Hence
\[
 \deg P_{13}=16(25+20+22)=1072.
\]
Its leading form is
$Q_0^{16}(2U_0V_0)^{16}(729U_0^2S_0)^8$, now with
$Q_0=a_0^4i_0^2c^4$ and $V_0=-a_0^2i_0^2c^4R_0$ from the
unretained-strong definitions. The explicit certificate is
\[
 [j^{480}w^{160}v^{192}h^{32}r^{64}c^{128}z^{16}]P_{13}
 =2^{1000}3^{696}\beta^{480}>0.
\]

\subsection{Literal arithmetic evaluation}
The complete DAG counts are 195, 148, 130, and 118 operations for
13, 14, 15, and 16 integers. They include the finalizer and all scalar
products. The product factors are evaluated after cancellation of their
constant terms, as displayed in Section~\ref{sec:integer}.
The five-root evaluator also shares terms in the first biquadratic
square when its two linear components coincide. These are exact ring
identities, not on-solution substitutions. The circuit files enumerate
every operation, and the verifier checks topology and liveness.


\section{A direct quartic from proof certificates}\label{sec:direct}
This branch represents the actual finite proof-certificate calculus directly.
Its integer coefficients are small literal numbers. The construction does
not encode a K215 computation and does not specialize a Jones universal
polynomial. It retains the existing C222/Horn/CRT compiler and improves a
repeated arithmetic component with a complete soundness and completeness
proof. The decisive component is the Diophantine graph of positive powers.

\subsection{An eleven-coordinate certificate for a positive power}
Let $n,k>0$ and $m\in\N$. Use eleven natural auxiliary coordinates
\[
 U,a,z,y,q,d,g,b,s,r,h.
\]
Define, without adding coordinates,
\begin{align*}
 W&=n+k,&Z&=z+1,&Y&=k+y,&T&=m+g+1,\\
 A&=n+a+1,&V&=Y^2(d+1),&B&=1+4Y(q+1),\\
 X&=Y(a+1)+m+Th,&S&=X+Us,&R&=k+4Yr.
\end{align*}
For this section only, put
$\mathcal P(A,X,Y)=X^2+Y^2-1-A^2Y^2$. The six equations are
\begin{align}
 \mathcal P(A,X,Y)&=0,&\mathcal P(A,U,V)&=0,&\mathcal P(B,S,R)&=0,
 \label{eq:powernorms}\\
 B&=A+Ub,&2An&=T+n^2+1,&\mathcal P(W+1,A,WZ)&=0.
 \label{eq:powerother}
\end{align}

\begin{theorem}[Positive-power graph]\label{thm:power}
For $n,k>0$ and $m\in\N$, equations
\eqref{eq:powernorms}--\eqref{eq:powerother} have a natural solution
in the eleven displayed auxiliaries exactly when $m=n^k$.
\end{theorem}

The underlying Pell--Matiyasevich criterion is classical and is available
as \texttt{Pell.matiyasevic} and \texttt{Pell.eq\_pow\_of\_pell}
in the pinned mathlib source\cite{mathlib}. The new point here is the
strictly positive specialization, the direction of the congruence
quotients, and exact elimination of coordinates.

\subsection{An intermediate certificate and soundness}
Before eliminating $X,S,R$, take the fourteen coordinates
\[
 A,X,U,S,R,z,y,q,d,g,b,s,r,h.
\]
Keep $W=n+k$, $Z=z+1$, $Y=k+y$, $T=m+g+1$, $V=Y^2(d+1)$,
and $B=1+4Y(q+1)$. In addition to the four Pell norms impose
\begin{align*}
 B&=A+Ub,&S&=X+Us,&R&=k+4Yr,\\
 X+Yn&=YA+m+Th,&2An&=T+n^2+1.
\end{align*}
The last norm is
\[
 A^2=1+W^3(W+2)Z^2.
\]
Since $W\ge2$ and $Z\ge1$, it gives $A>W\ge n$, in particular
$A\ge2$. All the standard criterion's size conditions now follow:
$Y\ge k>0$, $V\ge Y^2$, $B\ge5$, and $T>m$.
The directed equalities imply the required congruences modulo
$U,4Y,T$. The original positive-input criterion therefore gives $m=n^k$.
The inequality $A>n$ also ensures that the occurrence of $A-n$ in the
formal source is an ordinary nonnegative difference, rather than an
unintended truncated-subtraction expression.

\subsection{Why four signed quotients become natural}
This argument holds for any solution of the original positive-input
criterion, not just a selected completeness witness. Since
$V\ge Y^2\ge Y\ge1$, its first two Pell norms give $U\ge X$ and
$U\ge A$. Also $B\ge2$ and the third norm gives $S\ge1$.

If $0<c\le u$, a positive integer congruent to $c$ modulo $u$ must
have the form $c+uq$ with $q\ge0$. This includes $c=u$, when the least
positive residue-zero representative is $u$. Apply it to $(A,U)$ and
$(X,U)$ to obtain the first two natural quotients. Since
$0<k\le Y<4Y$, the congruence for $R$ similarly gives $R=k+4Yr$ with
$r\ge0$.

For the last quotient, the first norm implies
\[
 X^2-(A-1)^2Y^2=1+2(A-1)Y^2>0,
\]
hence $X>(A-1)Y$. Therefore
\[
 D=X+Yn-YA-m>(n-1)Y-m\ge-m>-T.
\]
The old power congruence makes $D$ a multiple of the positive $T$.
Every negative multiple is at most $-T$, so $D\ge0$ and $D=Th$ with
$h\in\N$. This proves the sign, rather than inferring it from the
positivity of some summands.

\subsection{Fixing the bound \texorpdfstring{$W$}{W} and proving completeness}
The standard completeness proof chooses $W=\max(n,k)$, but uses only
$W\ge n,k$ and $W>0$. Its explicit construction works with $W=n+k$.
Write
\[
 \chi_j(A)+\psi_j(A)\sqrt{A^2-1}
 =(A+\sqrt{A^2-1})^j.
\]
Choose
\[
 A=\chi_W(W+1),\qquad Z=\psi_W(W+1)/W.
\]
The recurrence modulo $W$ gives $W\mid\psi_W(W+1)$, and
$\psi_W(W+1)\ge W>0$, so $Z$ is a positive integer. Its norm gives
the last equation. Standard Pell growth yields
\[
 n^k\le n^W<(W+1)^W\le A.
\]
The strict comparison is between bases $n<W+1$, not between the first
two powers. Put $X=\chi_k(A)$, $Y=\psi_k(A)$, and
$T=2An-n^2-1$. Then $Y\ge k$ and, since $A\ge n+1$,
\[
 T-A=(2n-1)A-n^2-1\ge n^2+n-2\ge0.
\]
Hence $T\ge A>m=n^k$. Reduction in the quadratic ring modulo $T$,
with the substitution $\sqrt{A^2-1}=n-A$, gives
$X+Yn\equiv YA+n^k\pmod T$.
The standard Pell-index certificate supplies $U,V,S,R,B$ with the
remaining congruences and positive $Y^2$-multiple $V$.

The proved quotient signs and the nonnegative slacks
\[
 y=Y-k,\quad d=V/Y^2-1,\quad g=T-m-1,\quad z=Z-1,
 \quad q=(B-1)/(4Y)-1
\]
then give the fourteen-coordinate certificate. The last quotient is
nonnegative because $B>1$, $Y\ge1$, and $B\equiv1\pmod{4Y}$.

Finally $A>n$ permits the bijective substitution $A=n+a+1$.
The three directed equalities solve exactly as
\[
 X=Y(a+1)+m+Th,\qquad S=X+Us,\qquad R=k+4Yr.
\]
These are polynomials with nonnegative coefficients. They eliminate
three coordinates and their equations while preserving natural witnesses
in both directions. The eleven-coordinate theorem follows.

\subsection{The complete source and final quartics}
The retained compiler checks the literal C222 numerical proof rules,
uses a single finite-stream certificate and uniform CRT coding, and
expresses its binomial, factorial, and power operations by supplied
arithmetic graphs. There are exactly 55 power calls, 16 binomial calls,
and two factorial calls in the complete graph. Their input-positivity
dependencies are acyclic. The outer length and label bound, and each
encoded local-coordinate stream, are independent natural coordinates plus one.
Thus the carrier base and exponent are positive, and each range-binomial
call has positive upper and lower arguments. The factorial inputs are
the positive stream length and a sum with an added one. Within each
factorial graph, the radius base and exponent are the input plus two;
the subsequent binomial and power inputs are therefore positive. Finally,
the factorial output used as a base in the modulus-product graph is at
least one before that power call is evaluated. Induction in this macro
dependency order proves positivity at all 55 calls; the supplied symbolic
audit checks the actual complete call inventory.

These conventions preserve completeness. A refutation of target label
one has positive length, its label bound can be enlarged, and a CRT
representative can be made positive by adding a positive multiple of its
modulus. Thus the new graph is used only on its proved domain, without
excluding a genuine finite refutation. It is not asserted to represent
zero-base or zero-exponent cases.

One further quotient can be oriented globally. The local row polynomial
is a product of sums of squares plus additional squares, written as
$P=E-F$. Hence $E-F\ge0$ on every natural assignment. Its final
congruence modulo a positive $B$ can be written $E=F+Bq$, $q\in\N$,
instead of representing an arbitrary signed quotient by two naturals.

Every arithmetic wire is computed by a triangular addition or
multiplication of natural quantities. Eliminating a monic affine defining
row whose restoration has nonnegative coefficients is a bijection on
natural solution sets. The builder performs exactly those eliminations,
with a support cap that affects size but not correctness.

The delivered main polynomial is literally
\begin{equation}
 \boxed{D_{2473}(x_0,\ldots,x_{2472})
   =\sum_{i=1}^{2089}q_i(\boldsymbol{x})^2=0,
   \qquad\boldsymbol{x}\in\N^{2473}.}\label{eq:directquartic}
\end{equation}
All $q_i$ are the explicitly enumerated quadratic integer polynomials in
\texttt{direct/eq/equation.json}. The same polynomial is fully expanded
in \texttt{expanded\_polynomial.json}; thus the formula does not conceal
an unspecified polynomial or a compiler-existence assertion. For example,
the first residual is
\[
 x_{13}-1-2x_0-x_0^2-2x_0x_9-2x_9-x_9^2.
\]
The exact polynomial has 20,580 monomials, degree four, and coefficient
height 127,808,289. The alternative retaining affine wires has 4,860
natural variables, 4,476 quadratic rows, 14,719 monomials, and height
1,837. Every residual coefficient of that variant has absolute value
at most four.

\begin{theorem}[Direct consistency equation]
With the supplied literal source calculus and inherited ordinary
C222/Horn/CRT calibration,
\[
 \mathrm{ZFC}\vdash
 \bigl(\neg\exists\boldsymbol{x}\in\N^{2473}\;
 D_{2473}(\boldsymbol{x})=0\bigr)\longleftrightarrow\Con(\mathrm{ZFC}).
\]
The analogous assertion holds for the 4,860-variable variant.
\end{theorem}
\begin{proof}
Theorem~\ref{thm:power} and the final quotient argument preserve the
complete arithmetic graph's existential relation. The triangular affine
eliminations preserve natural solutions, and a sum of integer squares
is zero exactly when every residual is zero. The inherited ordinary
source calibration identifies the original graph with existence of a
finite C222, hence ZFC, refutation. Compose these equivalences and negate.
\end{proof}
The independent checker expands every residual square and matches every
coefficient, verifies all declared variables occur, and checks the 55
power-block census. It imports neither the author's builder nor a symbolic
algebra package. Its finite result supports the transcription; the
unbounded source calibration remains the stated mathematical dependency.


\section{Removing an unnecessary congruence restriction}\label{sec:freej}
This companion result extends the domain of an existing positive-variable
construction. Its numerical saving is only two operations; its new
argument proves that a formerly excluded congruence class cannot create
a spurious solution.

This section uses the separate, coherent compact44 tuple
$\beta_+=2^{590069}-1$, $K_0,M_C,M_F$, with $M_F$ already shifted.
It is not interchangeable with Section~\ref{sec:integer}'s source-34
tuple. The four exact payloads and their finite recipe are retained under
\texttt{positive/coefficients/}.

\subsection{The complete positive-variable polynomials}
For twelve strictly positive coordinates use
\[
 J,F,\alpha,t,T,X,Y,r,\eta,\zeta,n,\gamma.
\]
For thirteen use
\[
 J,F,\alpha,f,i,T,X,Y,r,\eta,\zeta,u,\gamma.
\]
All symbols in this subsection are local to these formulas. Define
\begin{align*}
 q&=\beta_+J+1,&C&=q-F-\alpha,&E&=XY,&k&=\eta+\zeta,\\
 c&=kY+\eta,&a&=(X+1)Y,&H&=4a+3,&\Delta&=(a+2)^2-1,\\
 \ell&=EkY,&M&=q^3(q-1),&W&=(q-1)k+1,&U&=q^2Y,\\
 B_t&=(qK_0+X)C-qF,&&&&&
\end{align*}
and
\begin{align*}
 Z&=M(\ell+k+r)+(q-1)Y-q^3WB_t,\\
 A_0&=Z^2-M^2[\ell(\ell+k)+1],\\
 R&=[(q-1)(q-F)+\alpha](q^2-1)+(M_C+qM_F)J,\\
 v&=k-R-1,\\
 D_*&=U(X+ac)+H(U\gamma+qE-v),\\
 A_1&=D_*^2-U^2(1+\Delta c^2).
\end{align*}
In both tails put $Q=\Delta(f^2-1)$ and $V=f^2(cT-R)-c$.
In the twelve-coordinate case $f=t+1$ is an abbreviation. Let
\[
 d=c^4,\quad h_0(z)=(1-Q)z^2+2(1-Q)Vz+V^2-1,
\]
\[
 b_n=2(1-Q)(n+V),\quad A_n=d\,h_0(n)+(1-Q)Q,\quad
 L=A_n^2-d\,b_n^2Q.
\]
The twelve-positive equation is
\[
 P_{12}^+=(1+A_0)(1+A_1)(1+L)-1=0.
\]
For thirteen positive coordinates set $I=ic^2$ and $y_*=V+u$, giving
\[
 P_{13}^+=(1+A_0)(1+A_1)
       [QV^2+(1-Q)y_*^2][1+I^2-Q]-1=0.
\]
Their complete DAGs have 98 and 88 operations and exact degrees 70 and 51.
The parents first compute $J=6j+1$ or $J=3j+1$; the new versions accept
positive $J$ directly and remove those two operations.

\subsection{Recovering divisibility despite a common factor three}
The product finalizers still recover their intended equations without
assuming $q$ even. Indeed $M=q^3(q-1)$ is always even;
$\Delta\equiv0$ or $3\pmod4$, so $A_1$ is never $2\pmod4$.
The remaining tail factors have the same exclusion of $-1$ as in the
ordinary parent proofs. Thus their unit argument applies on the enlarged
positive domain.

The unchanged first carrier proves $C>0$, $q^3\mid Y$, $q\ge16$,
and $3q<R<q^4-2$. The main carrier then gives a positive integer
$D=D_*/U$ with $D^2-\Delta c^2=1$ and
\[
 U\mid H(qE-v).
\]
The needed weaker coprimality statement is exact:
\begin{equation}
 \gcd(q,H)=\gcd(q^5,H)=\gcd(q,3)\in\{1,3\}.\label{eq:gcd3}
\end{equation}
For primes other than three this follows from $H\equiv3\pmod{q^3}$.
If $3\mid q$, then $27\mid Y$ and $H\equiv3\pmod{27}$, so its
$3$-adic valuation is exactly one. Since $q^5\mid U$, the displayed
divisibility still yields $q^4\mid(qE-v)$, and $q^4\mid qE$ gives
$q^4\mid v$.

The retained first/main Pell classification supplies indices $n_0,p$ with
\[
 k=2\psi_{2XY^2+1}(n_0),\quad c=\psi_{a+2}(p),\quad p>n_0,
 \quad 2n_0\equiv R+1\pmod{q^4}.
\]
Consequently $n_0\ge25$. The early rank bounds and the parent ordinary
strong and auxiliary arguments, now licensed by $q^4\mid v$, give
$p=R$ and then $R=2n_0-1$ exactly. These arguments do not use the
discarded residue restriction. In particular $U\mid v$, and the main
carrier gives
\begin{equation}
 q\mid3X,\qquad X\equiv2^R\pmod{H/g},\quad
 g=\gcd(q,H)\in\{1,3\}.\label{eq:weakprojection}
\end{equation}
The new task is to turn this weaker projection into equality before
assuming that $q$ divides $X$.

\subsection{Pell growth eliminates the exceptional branch}
Write $R=2r_0+1$, with $r_0=n_0-1\ge24$, and $P_0=2XY^2+1$,
$A=(X+1)Y+2$. Pell growth gives
\[
 \frac ck\ge\frac12\left(\frac{(2A-1)^2}{2P_0}\right)^{r_0},
 \qquad\frac ck<Y+1.
\]
The exact identity
\[
 X(2A-1)^2-2P_0(X+1)^2
 =12XY(X+1)+9X-2(X+1)^2>0
\]
for $X,Y\ge1$ improves this to
\begin{equation}
 Y+1>\frac12\left(\frac{(X+1)^2}{X}\right)^{r_0}.\label{eq:growth}
\end{equation}
For $X\ge2$, the base is at least $9/2$. The elementary inequality
$9^{16}>6\cdot8^{16}$ and $r_0\ge24$ imply $Y>2^R$.
Both $X$ and $2^R$ lie in $(0,a)$, whereas $H/g\ge H/3>a$.
Their congruence in \eqref{eq:weakprojection} therefore gives $X=2^R$.

For $X=1$, equation~\eqref{eq:growth} implies $Y\ge2^{R-2}$.
Then $H=8Y+3$ and $0<2^R-1<H/2$. Congruence modulo $H$ is impossible;
modulo $H/3$, the only possible positive multiple is $H/3$ itself.
It would give $8Y=3\cdot2^R-6$, impossible modulo eight for $R\ge3$.
Thus $X=1$ is excluded and again $X=2^R$.

Now $q\mid3\cdot2^R$ makes $q=2^u$ or $q=3\cdot2^u$.
The authentic Mersenne modulus eliminates the latter. If
$2^d-1\mid3\cdot2^u-1$, reduce $u$ modulo $d$ to $r\in\{0,\ldots,d-1\}$.
For $r\le d-2$, the number $3\cdot2^r$ lies strictly between one and
$2^d-1$. For $r=d-1$, its residue is $2^{d-1}+1$, also not one.
Here $d=590069\ge4$. Hence $q$ is a power of two and $q\mid X$.
The original source witnesses can now be restored exactly as in the
parent. Canonical refutation witnesses already have positive $J$, so
completeness is unchanged. This proves the same ordinary consistency
equivalence on the larger domain.

The exact degree certificates remain nonzero after replacing the old
leading coefficient $6\beta_+$ or $3\beta_+$ of $q$ by $\beta_+$.
The retained leading coefficients are $64\beta_+^{18}$ and
$-32\beta_+^{15}$. The circuit check confirms the counts
$54M+25A+19S=98$ and $48M+23A+17S=88$.


\section{A product certificate that computes its own bound}\label{sec:budget}
This separate theorem concerns a prospective alternative source. It is
not included in the measured cost claims for the delivered quartics.

Consider a binary Horn calculus on natural-number labels. The derivable
labels form the least set containing zero and closed under the permitted
rule instances $v\leftarrow(p,q)$. Thus zero is true and may be used as a
dummy premise; any other axiom is represented by a rule
$v\leftarrow(0,0)$. Label one is the target. For $N$ instances set
\[
 E=\prod_{i=0}^{N-1}(1+v_i)(1+p_i)(1+q_i),\qquad M=2E+2.
\]
Let $C_0=1$, $C_N=E$, and use the local recurrences
\begin{align}
 C_{i+1}&=C_i(1+v_i)(1+p_i)(1+q_i),\label{eq:budgetrec}\\
 S_{i+1}(1+Mv_i)&=S_i(1+Mp_i)(1+Mq_i),\label{eq:flowrec}
\end{align}
with $S_0=M+1$ and $S_N=1$.

\begin{theorem}[Product budget]
There exist $N$, permitted rule instances, and positive integer streams
satisfying these equations if and only if label one has a finite derivation.
\end{theorem}
\begin{proof}
Telescoping gives $A(M)=B(M)$ for
\[
 A(T)=\prod_i(1+v_iT),\qquad
 B(T)=(1+T)\prod_i(1+p_iT)(1+q_iT).
\]
The coefficients are nonnegative. Their sums are $A(1)\le E$ and
$B(1)\le2E<M$, so every coefficient is a base-$M$ digit. Uniqueness of
base expansion gives equality in $\Z[T]$. Unique factorization into
the normalized factors $1+aT$, for $a>0$, yields exact multiset balance:
\[
 \{v_i:v_i>0\}
 =\{1\}\mathbin{\uplus}\{p_i:p_i>0\}
          \mathbin{\uplus}\{q_i:q_i>0\}.
\]
If the target were outside the least set of derivable labels, count
occurrences outside that set. Every rule with an underivable conclusion
has at least one underivable premise, so the output count is at most
the premise count. The multiset balance makes it one greater, a contradiction.
This counting argument permits cyclic rule dependencies and repeated
instances; it does not require an acyclic ordering of the supplied list.

Conversely, expand the target backwards along a finite proof tree. First
compute $E$ from its labels, then $M$. Let $S_i$ be the product of $1+Ma$
over the currently demanded nonzero labels. These positive integers start
at $M+1$, end at one, and obey \eqref{eq:flowrec}. The running product
gives \eqref{eq:budgetrec}. There is no circular choice: the labels fix
$E$ and $M$ before the streams are assigned.
\end{proof}

This removes a common-multiple condition or an externally supplied power
bound from that flow interface. It does not itself produce a smaller
fixed-arity Diophantine equation. Both finite streams still require
unbounded sequence coding, with all variables and operations paid.
The latest direct compiler already uses an efficient single power bound;
adding the second stream is not assumed to improve its measured size.


\section{Verification, limitations, and further research}
\subsection{What the artifacts certify}
The archive includes literal polynomial residuals or arithmetic DAGs,
complete fixed numerical data, deterministic builders, independent checks,
and receipts. The root verifier runs the selected components and checks
the package manifest. Standard Python suffices for the integer and direct
components; the positive companion additionally uses SymPy.

The integer-family checks reconstruct all four native numerals from the
retained 943-tile recipe and repeat 69,230 convolution checks. They
authenticate all six payloads and 86 retained source members, check DAG
topology, liveness, and operation counts, and establish degree attainment
on a univariate specialization with the actual fixed coefficients.
For each of the four forms, 72 signed modular assignments compare the
literal evaluator with an independently multiplied conjugate product;
the zero-scale and zero-$c$ cases are included. These modular tests are
transcription diagnostics, not the unbounded proof.

A separate standard-library checker proves the rank-two, rank-three, and
rank-four cleared-norm identities by exact sparse Laurent-polynomial
arithmetic, matching all 6, 35, and 495 terms respectively. It checks
denominator cancellation and the unique radicand-free monomial used by
the modulo-four argument, independently of the circuit builder.

The direct checker reconstructs every coefficient of the full sum of
squares for both quartics. It checks variable indices, occurrence of every
declared coordinate, quadratic row degrees, expanded degree and height,
and the census of all 55 power graphs. A separate symbolic audit verifies
the eleven-auxiliary graph and its eliminated-coordinate identities.
Regeneration checks compare the emitted files byte for byte. The independent
export checker imports neither the author's builder nor an algebra library.

The positive-family checks authenticate the exact parent circuits and
coefficient bindings, verify the Pell ratio identity, check all circuit
inputs and gates, and check exact leading coefficients. Their finite
diagnostics include 896 residue cases, 6,336 gcd cases, and 15,132 Mersenne
cases; the written universal proof is required in addition to those tests.

Normal and optimized Python runs pass. The mathematical reviews of the
new integer mechanism and direct power graph were performed separately
from their construction; their scope and conclusions are retained. The
PDF was compiled from the delivered LaTeX and rendered for visual review.
No bounded search for polynomial zeros is offered as evidence of consistency.

\subsection{Limits of the conclusions}
No absolute optimum is proved for variables, degree, monomial count,
height, printed length, or arithmetic cost. The universal-representability
literature already permits fewer variables at high degree, including
Sun's ten-integer construction\cite{sun10} and the classical nine-natural
route. Recent formal work distinguishes representability, explicit degree
bounds, and executable coefficient recipes\cite{bayer}. Our comparisons
are between supplied fixed-instance polynomials at the pinned snapshot.

All new proofs are ordinary mathematical arguments. Independent review
within this work is AI-assisted, not external peer review. The native and
direct-calculus calibrations remain named mathematical dependencies;
this package does not claim an end-to-end formal proof or infer it from
finite computational checks. In particular the ordinary source-to-refutation
theorems and a sealed object-language ZFC proof artifact are different
levels of evidence.

\subsection{Further research questions}
\begin{enumerate}[leftmargin=*]
\item Can other retained-root choices improve the full integer degree--cost
frontier? Every candidate needs a separate residue proof excluding the
unwanted integer unit, as well as a literal cost and degree check.
\item Can the main and auxiliary roots be merged with positivity using
fewer coordinates without a full additional conjugate norm? Root separation
must be established before using conclusions that depend on squarehood.
\item Can the direct compiler exploit quotient orientation globally,
instead of only in the power graph and final CRT evaluation? Every
restored arithmetic coordinate must retain its intended natural value.
\item Can a different local proof calculus reduce the number of power
or sequence-coding calls? A replacement requires explicit transformations
of finite refutations in both directions; the current choice of C222 is
not asserted optimal.
\item Does the product-budget source become cheaper after both finite
streams and their common CRT scale are fully paid? Its elementary source
proof alone gives no numerical compiler bound.
\item Can the new ordinary arguments be checked against the sealed
object-language ZFC endpoint over each model's internal natural numbers?
This differs from validating the supplied polynomial data.
\item Can the native coefficient tuple be shortened while preserving the
same determinant and congruence calibration? Replacing it with a smaller
tuple from a different branch without proving that interface is invalid.
\end{enumerate}


\appendix
\section{Files, provenance, and incorporation}
The root directory contains this article, its complete LaTeX sources,
and three self-contained research components:
\begin{description}[style=nextline,leftmargin=2em]
\item[\texttt{integer/}]
The four unrestricted-integer equations, their division-free DAG builder,
ordinary proofs, exact six-coefficient payloads and finite generator,
retained parent proofs and licenses, verification receipts, and an
independent review of the new mathematical steps.
\item[\texttt{direct/}]
The 2,473- and 4,860-variable quartics, their fully expanded monomials,
literal C222 calculus and arithmetic compiler, the eleven-coordinate
power proof, pinned mathlib and repository calibration sources,
independent export checker, and an explicit integer companion.
\item[\texttt{positive/}]
The 98- and 88-operation positive-variable equations, the full free-$J$
proof, exact parent circuits, the separate compact44 coefficient tuple
and its generator, and normal/optimized verification receipts.
\end{description}

The direct integer companion uses
\[
 \sigma(x,y,z)=x^2+y^2+z(z+1),\qquad\sigma(\Z^3)=\N.
\]
Nonnegativity is immediate. Surjectivity follows by writing $4m+1$ as
three squares: exactly one is odd, giving
$4m+1=(2x)^2+(2y)^2+(2z+1)^2$.
Thus replacing every coordinate of $D_{2473}$ by an independent $\sigma$
gives an explicit degree-eight equation in $7,419$ unrestricted integers
with the same solvability status. The supplied residual JSON implements
this substitution. This companion demonstrates a direct integer route;
it does not compete with the smaller native integer forms.

\subsection*{Reproduction}
From the package root, with Python and SymPy installed:
\begin{lstlisting}
python3 verify.py
python3 -O verify.py
\end{lstlisting}
Each component also has an individual verifier described in its README
or proof file. To rebuild the article:
\begin{lstlisting}
cd report
pdflatex -interaction=nonstopmode -halt-on-error report.tex
pdflatex -interaction=nonstopmode -halt-on-error report.tex
\end{lstlisting}
The root manifest binds every delivered content file. Logs produced during
new runs are written outside those sealed content files or captured by
the root verifier. Generated Python cache directories are not part of the
deliverable. Source text and historical comparison snippets are explicitly
identified as inherited material; their old counts are not silently
relabelled as the new counts.

\subsection*{Incorporation}
This directory can be incorporated as a new proposal under
\path{Computability/DiophantineConsistency/con-zfc-diophantine-forms}.
Choose the numerical prefix at incorporation time. The deliverable does
not overwrite any repository source or claim to update its formal endpoint.
The snapshot, immediate source-archive hashes, and retained member hashes
provide a reproducible boundary for comparison. The archive's README
summarizes the main results and the exact proof dependencies.


\begin{thebibliography}{99}
\bibitem{repo} V. Reshetnikov and contributors, \emph{ProveIt}, snapshot
\repo. \url{https://github.com/VladimirReshetnikov/ProveIt/tree/a88aefd7ddab4771534bb049c67350e397c91b9a}.
Exact source files and archive hashes are recorded in the package.
\bibitem{sunpent} Z.-W. Sun, \emph{On universal sums of polygonal numbers},
Science China Mathematics \textbf{58} (2015), 1367--1396.
\url{https://arxiv.org/abs/0905.0635}.
\bibitem{oh} B.-K. Oh, \emph{Ternary universal sums of generalized
pentagonal numbers}, Journal of the Korean Mathematical Society
\textbf{48} (2011), 837--847. \url{https://arxiv.org/abs/0911.1181}.
\bibitem{mr} Y. Matiyasevich and J. Robinson, \emph{Reduction of an
arbitrary Diophantine equation to one in 13 unknowns}, Acta Arithmetica
\textbf{27} (1975), 521--553. \url{https://doi.org/10.4064/aa-27-1-521-553}.
\bibitem{sun10} Z.-W. Sun, \emph{Ten unknowns for Hilbert's tenth problem
over the integers}, arXiv:2609.01594v2 (2 September 2026).
\url{https://arxiv.org/abs/2609.01594}.
\bibitem{bayer} J. Bayer, M. David, M. Ha\ss ler, Y. Matiyasevich, and
D. Schleicher, \emph{Diophantine equations over $\Z$: universal bounds
and parallel formalization}, arXiv:2506.20909v3 (22 October 2026).
\url{https://arxiv.org/abs/2506.20909}.
\bibitem{mathlib} The mathlib contributors, \emph{PellMatiyasevic.lean},
commit \nolinkurl{81a5d257c8e410db227a6665ed08f64fea08e997}.
\url{https://github.com/leanprover-community/mathlib4/blob/81a5d257c8e410db227a6665ed08f64fea08e997/Mathlib/NumberTheory/PellMatiyasevic.lean}.
The exact source is retained with its license and hash.
\end{thebibliography}
\end{document}
